\documentclass[11pt,a4paper]{article}

\usepackage[utf8]{inputenc}
\usepackage{amsmath,amssymb,amsthm}
\usepackage{booktabs}
\usepackage{hyperref}
\usepackage[margin=1in]{geometry}
\usepackage{microtype}

\theoremstyle{plain}
\newtheorem{theorem}{Theorem}
\newtheorem{lemma}[theorem]{Lemma}
\newtheorem{corollary}[theorem]{Corollary}
\theoremstyle{definition}
\newtheorem{definition}[theorem]{Definition}
\theoremstyle{remark}
\newtheorem{remark}[theorem]{Remark}

\newcommand{\sem}[1]{[\![#1]\!]}
\newcommand{\Dom}{\mathcal{D}}
\newcommand{\Sub}{\mathcal{S}}

\title{Exact Mutation Adequacy for Access-Control Policies\\
is a Property of Guards, Not of a Language}

\author{%
Mohammad Thabet Hassan \and Ahmed Alameri \and Fahad Sadek \and
Omar Alraas \and Abdulrahman Rezki\\[2pt]
\normalsize Canadian University Dubai\\[2pt]
\normalsize Supervisor: Lobna AbuSerrieh
}

\date{}

\begin{document}
\maketitle

\begin{abstract}
Mutation adequacy is normally approximated. Deciding whether a surviving mutant is
semantically equivalent to the original is undecidable for programs, so tools either sample
the mutant space, leave equivalent mutants in the denominator, or ask a human to triage
survivors. We show that for access-control policies the question is decidable, and that
decidability is a property of a policy's \emph{guards} rather than of any particular policy
language: whenever each guard's induced partition of the subject space is determined by the
policy text, the policy-relative quotient is finite with computable witnesses, semantic
equivalence is a table comparison, and full coverage of the quotient kills every
non-equivalent mutant by construction. The result is independent of the evaluation rule,
which we verify under first-match, Cedar's forbid-overrides, Kyverno's any-rule-fails, and a
four-valued XACML combining algorithm.

We then ask how much published policy satisfies the condition, and answer it with a decision
procedure rather than an estimate: across three ecosystems' corpora, 753 of 805 policies lie
inside, with \emph{no policy left undetermined}, and what places a policy outside is one
construct per ecosystem rather than general expressiveness. We price what decidability buys:
of 38 mutants of our reference policy, 16 are provably equivalent, so a suite whose adequacy
is complete by construction is reported at 57.9\% by a pipeline that cannot decide
equivalence.

Two results are negative and we report them as such. The cheap coverage proxy for the
criterion carries almost no information on real corpora --- exactly zero bits on one --- and
the one predictive association in our data is confined to the policies that lie
\emph{outside} the fragment, where the theory does not apply. The witness construction that
underpins the positive results is machine-checked with an SMT solver, because an earlier
version of its proof was wrong.
\end{abstract}

\section{Introduction}

Mutation testing asks whether a change to a program would be noticed by its tests
\cite{demillo1978hints}. The measure is attractive and the arithmetic is not, because two of
its ingredients are hard. Deciding whether a surviving mutant is semantically equivalent to
the original --- and therefore should not count against the suite --- is undecidable in
general \cite{budd1982two}, and the mutant space is large enough that tools sample it
\cite{jia2011analysis}. Practical tools accept the resulting error, and a body of work
attacks the equivalent-mutant problem with partial methods
\cite{papadakis2015trivial,offutt2001mutation}.

This paper is about a setting where neither ingredient is hard. An access-control policy is
a short ordered list of rules over a small vocabulary of declared labels, and the questions
that are undecidable for programs are decidable here. That much is unsurprising. What we
think is worth stating is \emph{why}, because the reason turns out not to be a property of
the language we or anyone else designed.

\paragraph{Thesis.} Exact mutation adequacy for an access-control policy follows from one
condition on its guards --- that each guard's partition of the subject space is determined
by the policy text --- and from nothing else. It does not depend on the evaluation rule, on
the size of the decision domain, or on the language being anybody's in particular. Most
published policy meets the condition. The cheap coverage proxy that the theory suggests does
not measure it.

\paragraph{Contributions.}
\begin{enumerate}
  \item A condition on guard families (\emph{finite refinement}, Definition~\ref{def:refining})
    under which the policy-relative quotient is finite with computable witnesses
    (Theorem~\ref{thm:quotient}), equivalence is decidable (Theorem~\ref{thm:equivalence}),
    and coverage of the quotient decides the mutation score exactly
    (Theorem~\ref{thm:kill}, Corollary~\ref{cor:score}).
  \item Independence from the evaluation rule (Theorem~\ref{thm:combining}), verified under
    four rules drawn from deployed languages.
  \item A decision procedure for fragment membership, and its result on three published
    corpora: 526 of 548 XACML policies, 205 of 235 Kyverno policies and 22 of 22 Cedar
    policies inside, nothing undetermined (Section~\ref{sec:membership}).
  \item A measurement of what decidability is worth: the equivalent share of a real mutant
    set, and the resulting error of the two approximations a tool without it must use
    (Section~\ref{sec:price}).
  \item Two negative results: the coverage proxy carries near-zero information on real
    corpora, and the association between the proxy and fault detection is confined to
    policies outside the fragment (Section~\ref{sec:negative}).
  \item An SMT certificate for the witness construction (Section~\ref{sec:mechanised}),
    added because the construction's original proof was wrong in a way review did not catch.
\end{enumerate}

\paragraph{What is not new.} The coverage criterion we use --- witness every class of a
finite partition of the input space --- is all-combinations coverage over a category
partition in the sense of Ostrand and Balcer \cite{ostrand1988category}. We claim no novelty
for it, and Section~\ref{sec:related} says so at length. What we claim is the equivalence
between that criterion and mutation adequacy inside a characterised fragment, the
characterisation itself, and the measurements above.

\section{The policy language}
\label{sec:language}

We fix a language to make the development concrete. Section~\ref{sec:combining} shows that
almost none of it matters.

\begin{definition}[Subjects]
Let $T$ and $A$ be finite label domains (trust levels and action classes). A \emph{subject}
is a tuple
\[
  s = (t, a, c, i, j, G, K) \in \Sub
   = T \times A \times \Sigma^{*} \times \Sigma^{*} \times \Sigma^{*}
     \times \mathcal{P}_{\mathrm{fin}}(\Sigma^{*}) \times \mathcal{P}_{\mathrm{fin}}(\Sigma^{*}),
\]
where $c$ is a data classification, $i$ and $j$ are source and tool identifiers, $G$ is a
finite set of purpose tags and $K$ a finite set of capabilities. $\Sub$ is infinite: three
components range over arbitrary strings and two are arbitrary finite sets of them.
\end{definition}

\begin{definition}[Guards and policies]
A \emph{guard} is a predicate on $\Sub$ built from: membership of $t$ or $a$ in a named
finite set; membership of $c$, $i$ or $j$ in a named finite set; comparison of $c$ against a
named bound within a declared taxonomy; non-empty intersection of $G$ with a named finite
set; and existence of a $k \in K$ matching a named pattern, where a pattern is a literal or
a literal followed by a final wildcard. A \emph{rule} is a guard together with a decision in
a finite domain $\Dom$. A \emph{policy} $P$ is a finite ordered sequence of rules together
with a default $d_0 \in \Dom$.
\end{definition}

\begin{definition}[Semantics]
$\sem{P} : \Sub \to \Dom$ sends $s$ to the decision of the first rule whose guard holds at
$s$, or to $d_0$ if none does.
\end{definition}

\section{The quotient, and what is exact within it}
\label{sec:quotient}

\begin{definition}[Policy-relative equivalence]
For a policy $P$, let $s \sim_P s'$ hold when $s$ and $s'$ give the same truth value to
every guard occurring in $P$. Write $\Sub/{\sim_P}$ for the quotient and call its classes
\emph{cells}.
\end{definition}

\begin{definition}[Finite refinement]
\label{def:refining}
A guard family $\Gamma$ is \emph{finitely refining} if for every finite tuple
$g_1,\dots,g_n \in \Gamma$ the map $s \mapsto (g_1(s),\dots,g_n(s))$ has finite image, and
for every value in that image a subject realising it is computable from the syntax of
$g_1,\dots,g_n$.
\end{definition}

\begin{lemma}[Scalar components]
\label{lem:scalar}
A membership guard over a named finite set $N$ induces at most $|N| + 1$ classes on its
component, and a witness for each is either an element of $N$ or any string outside it.
Comparison against a bound in a declared taxonomy induces at most $|{\rm taxonomy}| + 1$.
\end{lemma}

\begin{proof}
The guard's value at $s$ depends only on which element of $N$ the component equals, or on
its being in none; that is $|N| + 1$ possibilities, each realised by the corresponding
element or by a fresh string, both computable from $N$. For the bound guards the value
depends only on the component's rank in the declared taxonomy, of which there are
$|{\rm taxonomy}|$, plus the unranked case.
\end{proof}

\begin{lemma}[Set-valued components]
\label{lem:setvalued}
Let $R_1,\dots,R_m$ be the named purpose sets occurring in $P$ and $p_1,\dots,p_m$ the
corresponding guards $p_\ell(s) \equiv G \cap R_\ell \neq \emptyset$. Then
$s \mapsto (p_1(s),\dots,p_m(s))$ has image of size at most $2^m$, and each achieved value
is realised by a subset of $\bigcup_\ell R_\ell$. The same holds for capability guards with
$m$ named patterns.
\end{lemma}

\begin{proof}
The image is a subset of $\{0,1\}^m$, giving the bound. If a value is achieved by some $G$,
it is achieved by $G \cap \bigcup_\ell R_\ell$, since each $p_\ell$ consults only that
intersection; this set is finite and computable. For capabilities, if a value is achieved by
some $K$ then it is achieved by choosing, for each pattern marked true, one member of $K$
matching it, since dropping capabilities cannot make a pattern true that was false.
\end{proof}

\begin{remark}[The bound is not a count]
\label{rem:notacount}
Lemma~\ref{lem:setvalued} bounds the image by $2^m$ and does \emph{not} claim every element
of $\{0,1\}^m$ is achieved, and an earlier version of this development did claim exactly
that. It is false in both set-valued components, for different reasons. Every capability
matching \texttt{net.http} also matches \texttt{net.*}, so on a policy naming both patterns
the value $(\mathrm{false},\mathrm{true})$ is achieved by no subject at all. And a policy
whose only purpose guard is intersection with $\{a,b\}$ cannot distinguish $\{a\}$ from
$\{b\}$ from $\{a,b\}$, so three subsets realise one value. Enumerating $\{0,1\}^m$
therefore both invents classes no subject occupies and counts classes twice.
Section~\ref{sec:mechanised} machine-checks the corrected construction.
\end{remark}

\begin{theorem}[Finite quotient]
\label{thm:quotient}
For every policy $P$ over the language of Section~\ref{sec:language}, $\Sub/{\sim_P}$ is
finite, $\sem{P}$ is constant on each cell, and a witness for each cell is computable from
$P$. Writing $X_P, I_P, J_P, G_P, K_P$ for the classifications, source identifiers, tool
identifiers, purpose tags and capability patterns $P$ names,
\[
  |\Sub/{\sim_P}| \;\le\;
  |T| \cdot |A| \cdot (|X_P| + |{\rm taxonomy}| + 1) \cdot (|I_P| + 2) \cdot (|J_P| + 2)
  \cdot 2^{|G_P|} \cdot 2^{|K_P|}.
\]
\end{theorem}

\begin{proof}
$\sim_P$ is the kernel of the map sending $s$ to the tuple of truth values of the guards of
$P$. By Lemmas~\ref{lem:scalar} and~\ref{lem:setvalued} each component of that tuple takes
finitely many values with computable witnesses, and the displayed product bounds the number
of combinations. $\sem{P}$ consults only those truth values, so it is constant on each
class. A witness for a class is obtained by taking, componentwise, the witnesses the two
lemmas supply.
\end{proof}

\begin{theorem}[Decidable equivalence]
\label{thm:equivalence}
For policies $P$ and $Q$ over the language, $\sem{P} = \sem{Q}$ is decidable.
\end{theorem}

\begin{proof}
Both are constant on the classes of $\sim_P \cap \sim_Q$, which is the kernel of the map
given by the guards of $P$ and $Q$ together and is finite with computable witnesses by
Theorem~\ref{thm:quotient} applied to the concatenation of the two rule sequences. Compare
the two decisions on one witness per class.
\end{proof}

\begin{definition}[Suites, witnessing, killing]
A \emph{suite} $\Sigma$ is a finite set of pairs $(s,d) \in \Sub \times \Dom$. Its
\emph{witnessed set} is $W(\Sigma) = \{s : (s,d) \in \Sigma\}$, read up to $\sim_P$.
$\Sigma$ is \emph{consistent} with $P$ when $d = \sem{P}(s)$ for every pair. $\Sigma$
\emph{kills} $M$ when $\sem{M}(s) \neq d$ for some pair. Write
$\Delta(P,M) = \{s : \sem{P}(s) \neq \sem{M}(s)\}$.
\end{definition}

\begin{theorem}[Exact kill criterion]
\label{thm:kill}
For consistent $\Sigma$: $\Sigma$ kills $M$ if and only if
$\Delta(P,M) \cap W(\Sigma) \neq \emptyset$.
\end{theorem}

\begin{proof}
($\Leftarrow$) Let $s$ lie in both. There is a pair $(s,d) \in \Sigma$, and consistency
gives $d = \sem{P}(s)$. As $s \in \Delta(P,M)$ we get $\sem{M}(s) \neq \sem{P}(s) = d$, so
that pair fails. ($\Rightarrow$) If $\Sigma$ kills $M$ via $(s,d)$ then
$\sem{M}(s) \neq d = \sem{P}(s)$, so $s \in \Delta(P,M)$, and $s \in W(\Sigma)$ by
definition.
\end{proof}

\begin{corollary}[Coverage decides the score]
\label{cor:score}
If $W(\Sigma) = \Sub/{\sim_P}$ then $\Sigma$ kills every mutant not equivalent to $P$, for
any mutation operator set whatsoever.
\end{corollary}

\begin{proof}
A non-equivalent mutant has $\Delta(P,M) \neq \emptyset$ by Theorem~\ref{thm:equivalence},
and it meets $W(\Sigma)$.
\end{proof}

\begin{remark}[The converse, and what it quantifies over]
If $W(\Sigma) \neq \Sub/{\sim_P}$, pick a cell $s$ outside it and $d \neq \sem{P}(s)$. The
function agreeing with $\sem{P}$ except at $s$, where it returns $d$, is non-equivalent and
survives $\Sigma$. So over \emph{semantic} mutants --- arbitrary functions
$\Sub/{\sim_P} \to \Dom$ --- full coverage is necessary as well as sufficient. Over the
policies a given operator set generates it need not be, and a syntactic operator set can
therefore report a complete score against an incomplete suite. This is why an
implementation should report witnessed cells beside the score.
\end{remark}

\section{The evaluation rule does not matter}
\label{sec:combining}

Nothing above used first-match evaluation. Theorem~\ref{thm:kill} and
Corollary~\ref{cor:score} consult only the finite class set, the decision on each class, and
a witness per class.

\begin{theorem}[Combining-rule independence]
\label{thm:combining}
Let $P = (\langle g_i \rangle_{i \le n}, c)$ where the $g_i$ are drawn from a finitely
refining family, each $g_i$ reports a value in a finite set $V$, and
$\sem{P}(s) = c(g_1(s),\dots,g_n(s))$ for some $c : V^n \to \Dom$. Then $\Sub/{\sim_P}$ is
finite with at most $|V|^n$ classes, $\sem{P}$ is constant on each, witnesses are computable,
and Theorems~\ref{thm:equivalence} and~\ref{thm:kill} and Corollary~\ref{cor:score} hold
verbatim.
\end{theorem}

\begin{proof}
$\sim_P$ is the kernel of $s \mapsto (g_i(s))_i$, finite by Definition~\ref{def:refining};
$\sem{P} = c \circ (g_i)_i$ is therefore constant on each class. Every later argument uses
only those three facts, and none inspects $c$.
\end{proof}

$V$ need not be two-valued: an XACML guard also reports \emph{indeterminate} when a required
attribute is absent, which is what produces the \texttt{Indeterminate} decision. Table
\ref{tab:combining} instantiates $c$ for five deployed languages; all five are functions of
the guard outcome vector, so all five satisfy the hypothesis.

\begin{table}[t]
\centering
\begin{tabular}{llc}
\toprule
Language & $c$ on the guard outcome vector & $|\Dom|$ \\
\midrule
This paper & decision of the first true guard, else the default & 3 \\
Cedar \cite{cedar2024} & \texttt{Deny} if any \texttt{forbid}; else \texttt{Allow} if any \texttt{permit}; else \texttt{Deny} & 2 \\
Kyverno \cite{kyverno} & \texttt{fail} if a rule's pattern is violated; \texttt{pass} if all match & 3 \\
XACML \cite{xacml30} & the chosen combining algorithm & 4 \\
Rego (Gatekeeper) \cite{gatekeeper} & violation set non-empty iff any \texttt{violation} guard holds & 2 \\
\bottomrule
\end{tabular}
\caption{The evaluation rule is where deployed languages differ, and Theorem
\ref{thm:combining} shows it is the part the results do not use. We verify the theorem
computationally under the first four.}
\label{tab:combining}
\end{table}

\paragraph{The boundary is the guards.} What excludes a language is
Definition~\ref{def:refining}, not $c$. Guards whose partition is fixed by the policy text
are inside: equality or membership against literals, a final-wildcard pattern, a comparison
against a literal threshold, bounds in a declared taxonomy. Guards whose partition depends
on data the policy does not contain are outside: transitive membership over an entity store
of unbounded depth, a pattern taken from the input, an arbitrary built-in over a whole
document, or a reading of the clock.

\section{How much published policy is inside}
\label{sec:membership}

Membership is decidable by inspection of guards, so we decided it, on the corpora of three
deployed languages at pinned commits. The procedure refuses rather than guesses: a policy is
reported inside only when every function it names belongs to a family whose partition is
fixed by a literal in the policy, outside when it reads data the policy does not contain,
and undetermined otherwise.

We report two scopes, because they answer different questions. The \emph{joined} scope is
the set of policies a suite study could also score, and it is the scope the predictive
experiment of Section~\ref{sec:negative} stratifies; it is small because measuring decision
coverage needs a suite with more than one case. The \emph{whole-corpus} scope is every
policy the corpus contains, which is the right scope for asking how much of an ecosystem the
fragment covers, since membership is decided from policy text and needs no suite at all.

\begin{table}[t]
\centering
\begin{tabular}{lrrrrr}
\toprule
Ecosystem & Policies & Inside & Outside & Undetermined & Share inside \\
\midrule
\multicolumn{6}{l}{\emph{Whole corpus}} \\
XACML \cite{xacml30} & 548 & 526 & 22 & 0 & 96.0\% \\
Kyverno \cite{kyverno} & 235 & 205 & 30 & 0 & 87.2\% \\
Cedar \cite{cedar2024} & 22 & 22 & 0 & 0 & 100.0\% \\
\midrule
\textbf{Total} & \textbf{805} & \textbf{753} & \textbf{52} & \textbf{0} & \textbf{93.5\%} \\
\midrule
\multicolumn{6}{l}{\emph{Joined to a suite study} --- the scope Section~\ref{sec:negative} stratifies} \\
XACML & 21 & 15 & 6 & 0 & 71.4\% \\
Kyverno & 49 & 41 & 8 & 0 & 83.7\% \\
Cedar & 22 & 22 & 0 & 0 & 100.0\% \\
\bottomrule
\end{tabular}
\caption{Fragment membership. Nothing is undetermined in any ecosystem at either scope, so
each share is a verdict on the whole corpus rather than on the part the procedure happened
to understand. The joined scope is a subset of the whole-corpus scope and every verdict
agrees between them.}
\label{tab:membership}
\end{table}

What places a policy outside is one construct per ecosystem, not general expressiveness. In
XACML it is an \texttt{AttributeSelector} or an XPath function, either of which puts an
arbitrary XML document in the subject, or \texttt{access-permitted}, which re-enters the
decision point on a request the policy constructs. In Kyverno it is a \texttt{context} entry
that queries the cluster or a registry, a CEL call on the \texttt{resource} receiver that
reaches the API server, or \texttt{now()}. Cedar has no such construct: it cannot read data
the request does not carry, which is consistent with its being designed for automated
analysis \cite{cedar2024}.

\paragraph{Membership is a property of families, not of datatypes.} XACML names a function
\texttt{<type>-<family>}, and the family decides membership: \texttt{integer-greater-than}
and \texttt{time-greater-than} both compare a designator against a literal and both split
the subject space at that literal. Deciding by enumerated name rather than by family
under-reports the fragment, and measurably so --- the conformance corpus exercises 182
function names our first, name-enumerated allowlist did not carry, and all but one of them
belong to a family that allowlist already accepted for some other datatype. The exception is
\texttt{xpath-node-count}, which reads the request's \texttt{Content}, and it is outside for
the same reason an \texttt{AttributeSelector} is.

\paragraph{A policy may state no predicate at all.} 13 policies in the XACML corpus contain
no \texttt{Match}, \texttt{Condition}, \texttt{Apply} or \texttt{VariableDefinition}
element. These are not unjudged: a policy that states no predicate induces a quotient with a
single class, which is the fragment's strongest case rather than a gap in it. Distinguishing
them from a policy whose guards our scan failed to read matters for soundness, so the
procedure decides by the presence of guard-bearing elements and not by whether its function
scan came back empty.

\paragraph{Limits.} These are the projects' own test and conformance corpora, not deployed
policy. Membership is a property of guards, so a policy inside the fragment gets exact
adequacy for its own decision structure, and its language remains undecidable in general
(Section~\ref{sec:limits}).

\section{What decidability is worth}
\label{sec:price}

Corollary~\ref{cor:score} says an exact score is available. It does not say what the
alternative costs, so we measured that where the exact answer is known.

Our reference policy induces 12 cells and its operator set generates 38 mutants, of which
\textbf{16 are provably equivalent} by Theorem~\ref{thm:equivalence} --- 42.1\%. A tool that
cannot decide equivalence must either ask a human to triage every survivor or leave the
equivalent mutants in the denominator. Table~\ref{tab:price} gives the second.

\begin{table}[t]
\centering
\begin{tabular}{lrrr}
\toprule
Suite & Exact score & Without equivalence detection & Understated by \\
\midrule
default & 63.6\% & 36.8\% & 26.8 pt \\
adversarial & 36.4\% & 21.1\% & 15.3 pt \\
coverage-matrix & \textbf{100.0\%} & \textbf{57.9\%} & \textbf{42.1 pt} \\
\bottomrule
\end{tabular}
\caption{The error has one sign: undetected equivalents inflate the denominator, so a score
is always understated. The last row witnesses every cell, so its adequacy is complete by
Corollary~\ref{cor:score}, and is reported at 57.9\%.}
\label{tab:price}
\end{table}

Sampling costs separately. Because the true kill set is known, the error is computed rather
than simulated. Against the 63.6\% suite, a sample of 4 of the 22 non-equivalent mutants is
off by 0.19 on average and by more than ten points in \emph{every} sample; at 12 it exceeds
ten points in 38\% of samples; only at 20 of 22 does it reliably come within ten.

\section{The proxy does not measure the criterion}
\label{sec:negative}

Corollary~\ref{cor:score} requires coverage of the quotient, which is cheap to check for one
policy and needs the policy parsed. A tool wanting a language-independent proxy would
instead count \emph{decisions witnessed}: does a suite ever elicit each element of $\Dom$?
We measured that proxy on four ecosystems and it does not carry the criterion.

\paragraph{It is nearly information-free.} A threshold partitions subjects into flagged and
unflagged, and if it flags almost none or almost all it distinguishes little whatever its
standing in a proof. Measuring the binary entropy of the flagged share: the natural
threshold --- witness all $|\Dom|$ decisions --- is the most informative available in 2 of 5
ecosystem/domain pairs and carries under 0.3 bits in 4 of 5. On Kyverno mutate rules it
flags all 88 subjects and carries \textbf{0.000 bits} --- exactly none. The informative threshold is
interior: on XACML, three of four decisions splits the corpus 57/43 and carries 0.985 bits,
against 0.276 for all four.

\paragraph{Its one association is outside the fragment.} On the Kyverno corpus, policies
whose suites witness a single decision kill fewer mutants than those witnessing more: a
difference in means of 0.149 at $p = 0.043$ one-sided over 49 policies. Stratifying by
fragment membership dissolves it. Inside the fragment, over the larger arm of 41 policies,
the difference is 0.066 at $p = 0.272$; outside it, over 8, it is 0.498 at $p = 0.018$
exactly. The pooled figure is therefore not evidence about the criterion where the criterion
is exact. We also tested whether the relation is graded --- whether witnessing more decisions
predicts detecting more faults --- and it is not, at $p = 0.157$ by a Jonckheere--Terpstra
permutation test.

\paragraph{What this does and does not say.} It does not refute Corollary~\ref{cor:score},
which requires cell coverage; the stratified flag records only whether more than one decision
was witnessed, which is a weak proxy for that. Nor is it evidence that the proxy works
outside the fragment: all eight of those policies are image or registry policies, and
image-policy suites being weak explains their scores without reference to coverage.

\section{The witness construction, machine-checked}
\label{sec:mechanised}

Remark~\ref{rem:notacount} records a false step in an earlier version of
Lemma~\ref{lem:setvalued}. It survived review, so we do not rely on review for it. We encode
the capability matching rule in SMT --- a final-wildcard pattern as a string prefix, a
literal as equality --- and ask Z3 \cite{z3}, for every candidate value in $\{0,1\}^m$,
whether any capability set realises it. The construction must produce exactly the achievable
ones: no unachievable value, or it invents a class, and none missing, or the quotient is
incomplete and Corollary~\ref{cor:score}'s premise is unreachable. Solver and construction
agree on 8 of 8 pattern sets, chosen to cover nested, disjoint, literal, chained and mixed
patterns. Two nested patterns give 4 candidates and 3 achievable; a chain of three
gives 8 and 5.

\section{Where the fragment ends}
\label{sec:limits}

\begin{theorem}
\label{thm:undecidable}
If guards may contain arbitrary computable predicates, policy equivalence is undecidable.
\end{theorem}

\begin{proof}
Reduce from the emptiness problem for the domain of a partial computable function, which is
undecidable by Rice's theorem \cite{rice1953classes}. Given such a function $f$, let $P$
have the single rule ``if $f$ halts on $s$ then \texttt{allow}'' with default \texttt{deny},
and let $Q$ be the policy with no rules and default \texttt{deny}. Then
$\sem{P} = \sem{Q}$ iff $f$ halts nowhere.
\end{proof}

So the condition of Definition~\ref{def:refining} is doing real work, and
Section~\ref{sec:membership} measures how often it is met rather than assuming it.

\section{Related work}
\label{sec:related}

\paragraph{Mutation testing and the equivalent-mutant problem.} Mutation analysis dates to
DeMillo, Lipton and Sayward \cite{demillo1978hints}; Jia and Harman survey its development
\cite{jia2011analysis}. Equivalence of mutants is undecidable in general
\cite{budd1982two}, and the literature accordingly pursues partial methods: trivial compiler
equivalence detects a substantial share cheaply \cite{papadakis2015trivial}, and Offutt and
Untch review the cost-reduction techniques including mutant sampling
\cite{offutt2001mutation}. Our contribution is not a better partial method. It is the
observation that a policy language whose guards satisfy
Definition~\ref{def:refining} makes the problem decidable outright, and the measurement in
Section~\ref{sec:price} of what that is worth.

\paragraph{Category-partition testing.} The criterion we use --- witness every class of a
finite partition of the input space --- is all-combinations coverage over a category
partition \cite{ostrand1988category}. We claim no novelty for the criterion. Our claim is
Corollary~\ref{cor:score}: inside the fragment it is provably \emph{equivalent} to mutation
adequacy, which turns an adequacy question with an undecidable component into a coverage
question costing one pass over the suite.

\paragraph{A naming collision.} ``Decision coverage'' in ISTQB usage and in DO-178C
\cite{do178c} means branch coverage: a structural criterion over program text. What we
measure is a criterion over a policy's \emph{output}. The two are not comparable, and a
suite can satisfy the structural one completely while failing the output one; we exhibit
such a suite. Any write-up must define the term on first use, and we have avoided it in this
paper's title and abstract for that reason.

\paragraph{Policy analysis.} Cedar is designed to admit automated reasoning and ships an
SMT-based analysis tool \cite{cedar2024}; Table~\ref{tab:membership} finding all of its
policies inside the fragment is consistent with that design rather than a surprise. XACML
\cite{xacml30} standardises the four-valued domain and the combining algorithms of
Table~\ref{tab:combining}. Kyverno \cite{kyverno} and Gatekeeper \cite{gatekeeper} are the
Kubernetes admission engines whose corpora we measure.

\section{Threats to validity}

\paragraph{Construct.} Our exactness claim is about the quotient. An implementation that
enumerates the product of per-component witness spaces is enumerating a \emph{refinement} of
the quotient: distinct product cells can answer every guard identically, and then they are
one class. Soundness carries over, since $\sem{P}$ is constant on any refinement of a
partition it is constant on, so Theorem~\ref{thm:kill} and Corollary~\ref{cor:score} hold as
stated. The policy-level reading of the converse does not: where two cells belong to one
class, no policy the language can express differs at one and not the other.

\paragraph{Internal.} The figures in Sections~\ref{sec:price} and~\ref{sec:negative} come
from one reference policy and four corpora, and the operator set is syntactic, so scores
measure suites against those edits rather than against real faults. Section
\ref{sec:mechanised} exists because one proof step was wrong; we do not claim the others are
machine-checked.

\paragraph{External.} The corpora of Table~\ref{tab:membership} are the projects' own test
and conformance directories, not deployed policy, and three ecosystems chosen by us are not a
sample of the field. The whole-corpus scope is dominated by XACML, and that corpus is a
conformance suite: it is written to exercise the specification's function library, so it
over-represents unusual functions and under-represents whatever policy shapes are common in
production. Reading its 96.0\% as the share of deployed XACML inside the fragment would be a
mistake. The Kyverno predictive experiment has nine policies in its blind arm, and its operator
set was extended between two runs --- a researcher degree of freedom we report rather than
hide.

\paragraph{Statistical.} The stratified analysis in Section~\ref{sec:negative} is a
subgroup analysis of an already marginal effect, and the out-of-fragment stratum is
confounded: all eight of its policies are image or registry policies, the group whose guards
read a registry, so weak image-policy suites explain their low scores without reference to
coverage.

\section{Conclusion}

Exact mutation adequacy for access-control policy does not require a special language. It
requires that each guard's partition of the subject space be determined by the policy text,
and that condition is met by 753 of the 805 published policies we examined across three
ecosystems, with none left undetermined. Where it holds, the equivalent-mutant problem is a table
comparison and coverage of the quotient decides the score --- worth 42 percentage points on a
suite whose adequacy is in fact complete. Where a cheap language-independent proxy is used
instead, it carries almost no information, and the one association we found between that
proxy and fault detection lies outside the fragment the theory covers. We report the second
half as plainly as the first.

\bibliographystyle{plain}
\bibliography{refs}

\end{document}
